\documentclass[11pt,a4paper]{article}

% ------------------------------------------------------------
% Paquetes
% ------------------------------------------------------------
\usepackage[utf8]{inputenc}
\usepackage[T1]{fontenc}
\usepackage[spanish,es-noshorthands,es-tabla]{babel}
\decimalpoint
\usepackage{lmodern}
\usepackage{amsmath,amssymb}
\usepackage{graphicx}
\usepackage{booktabs}
\usepackage{float}
\usepackage[margin=2.5cm]{geometry}
\usepackage[round,authoryear]{natbib}
\usepackage{xcolor}
\usepackage[colorlinks=true,linkcolor=blue!50!black,citecolor=blue!50!black,urlcolor=blue!50!black]{hyperref}
\usepackage{enumitem}
\setlist{itemsep=2pt,topsep=4pt}

\graphicspath{{figuras/}}

\newcommand{\porb}{P_{\rm orb}}
\newcommand{\tx}{\texttt{tess\_extractor}}
\newcommand{\Dum}{D_2/D_1}

\title{Profundidades de eclipse entre sectores de TESS en Orión:\\
¿variación real de las estrellas o efecto sistemático?}
\author{Notas de procedimiento}
\date{\today}

\begin{document}
\maketitle

\begin{abstract}
La prueba de consistencia de \texttt{1b\_profundidades\_eclipse\_v2.ipynb} marcó 13 de los 21
sistemas porque la profundidad de sus eclipses cambia de un sector de TESS a otro. Estas notas
explican cómo se separó la parte real (cambio del flujo de las estrellas) de la parte sistemática
(fotometría, cadencia y efeméride). Se compararon cinco fotometrías independientes de los mismos
sectores (\tx, TGLC, QLP, CDIPS y PDCSAP), se midió cada eclipse por separado en su propio mínimo y
se llevaron todas las curvas a una cadencia común de 30~min. Los resultados son tres. (1)~Las
curvas de \tx\ están diluidas: sus eclipses son entre 15 y 50~\% más someros que en las
fotometrías que restan las estrellas vecinas, como advierten los propios autores de la herramienta.
(2)~Buena parte de las diferencias entre sectores eran cadencia (30~min en S5/S6) y efeméride
imprecisa. (3)~En cuatro sistemas el cociente $\Dum$, que no depende de la dilución, cambia en el
mismo sentido en todas las fotometrías: ahí hay una variación real.
\end{abstract}

\tableofcontents

% ============================================================
\section{Resumen de resultados}
% ============================================================

\begin{table}[H]
\centering
\caption{Veredicto por sistema. ``$D_1$ entre sectores'' se refiere a la profundidad del primario
en un mismo método fotométrico; ``$\Dum$'' al cociente de profundidades, que es independiente de la
dilución (sección~\ref{sec:dilucion}). Todo a cadencia común de 30~min.}
\label{tab:resumen}
\small
\begin{tabular}{llll}
\toprule
Sistema & TIC & $D_1$ entre sectores & $\Dum$ \\
\midrule
1  & 176775425 & estable; S98 de \tx\ $+7$~\% (sistemático) & \textbf{real}: sube de 0.54 a 0.62 \\
3  & 11412481  & consistente & estable \\
4  & 50897998  & S98 $-5$~\% en \tx\ y PDCSAP (dilución común) & estable \\
5  & 249074181 & sin fotometría independiente & sin comparación \\
6  & 72921630  & sin fotometría independiente & sin comparación \\
7  & 264483228 & consistente & leve aumento (0.068 a 0.074) \\
8  & 11300044  & consistente & sin comparación \\
9  & 427353034 & S98 de \tx\ $-24$~\% (sistemático) & \textbf{real}: baja de 0.39 a 0.36 \\
10 & 72829000  & consistente & sin comparación \\
11 & 138829403 & S32 de \tx\ $-7$~\% (sistemático) & estable \\
12 & 264741359 & estable ($\pm1.5$~\%) & \textbf{real}: baja de 0.387 a 0.342 \\
16 & 427395536 & consistente & posible aumento (ruidoso) \\
18 & 427347964 & \textbf{real}: S32 $+25$~\%, S98 $-35$~\% & \textbf{real}: sube en S98 \\
\bottomrule
\end{tabular}
\end{table}

En pocas palabras:
\begin{enumerate}
\item Las curvas de \tx\ son fotometría de apertura sin corrección de contaminación. En casi todos
  los sistemas tienen entre 15 y 50~\% de flujo ajeno a la binaria (tabla~\ref{tab:dilucion}). Esto
  achica las profundidades y, en starry, los radios de la secundaria.
\item La fracción de flujo ajeno cambia de sector a sector. Por eso la profundidad de \tx\ cambia sin
  que cambie la estrella (sistemas 1, 9 y 11).
\item La cadencia de 30~min de S5/S6 achata los eclipses cortos entre 5 y 15~\% (figura~\ref{fig:cadencia}).
  Como S5/S6 son también los sectores más antiguos, esto parecía una tendencia temporal.
\item En los sistemas 1, 9, 12 y 18, $\Dum$ cambia en el mismo sentido en todas las fotometrías y a la
  misma cadencia. Eso no lo puede producir la dilución: es un cambio real del brillo relativo de las
  estrellas (manchas) o de la geometría.
\end{enumerate}

% ============================================================
\section{El problema}
% ============================================================
La profundidad de un eclipse es una propiedad geométrica del sistema: depende de los radios, de la
inclinación y del cociente de brillos superficiales. Salvo que alguna de esas cosas cambie, debería
ser la misma en todos los sectores. La prueba $\chi^2$ de la sección~7 del notebook 1b encontró que
no es así en 13 sistemas (figura~\ref{fig:plegadas1b}). Hay tres familias de explicaciones:
\begin{itemize}
\item \textbf{reales}: manchas que cambian el brillo de una de las estrellas, una tercera estrella
  variable en la apertura, o un cambio de inclinación por una tercera estrella \citep{Baum2025};
\item \textbf{fotométricas}: distinta contaminación por estrellas vecinas o por la nebulosa en cada
  sector;
\item \textbf{de la medición}: cadencia distinta en cada sector y efeméride imprecisa.
\end{itemize}

\begin{figure}[H]
\centering
\includegraphics[width=0.85\textwidth]{fig1_plegadas_1b.png}
\caption{Curvas plegadas por sector de \tx, tal como salieron del notebook 1b. Cada color es un
sector; la línea discontinua marca $1-D$ de cada sector.}
\label{fig:plegadas1b}
\end{figure}

% ============================================================
\section{Fotometrías}
\label{sec:fotometrias}
% ============================================================

\subsection{\tx\ (TESSExtractor)}
Las curvas de TESS \citep{Ricker2015} de \texttt{Light\_Curve\_Files/} vienen de TESSExtractor \citep{Serna2026}, la herramienta
que se usó para la rotación de estrellas T~Tauri en el complejo de Orión \citep{Serna2021}. Según su
descripción:
\begin{itemize}
\item descarga un corte de $10\times10$~píxeles de las imágenes completas (FFI) con TESScut;
\item hace fotometría de apertura con un radio que depende de la magnitud (1.5 a 3~píxeles);
\item estima el fondo con la \emph{moda} de los píxeles de un anillo alrededor de la estrella;
\item corrige sistemáticos con vectores de tendencia común (CBV): en una escala (\texttt{Flux\_CBV\_SS})
  y en varias escalas más picos (\texttt{Flux\_CBV\_MS}).
\end{itemize}
\textbf{No corrige la dilución}. Solo calcula una métrica de contaminación con las vecinas de Gaia y
recomienda tratar con cautela los objetos con más de 15~\%. Los mismos autores señalan que la
fotometría de apertura ``inevitablemente incluye flujo contaminante'' y da eclipses más someros que la
fotometría PSF de TGLC \citep{Serna2026}.

\subsection{TGLC}
TESS--Gaia Light Curve \citep{HanBrandt2023} modela la escena completa: con las posiciones y
magnitudes de Gaia DR3 ajusta una PSF efectiva (ePSF) en cada imagen, resta el flujo de todas las
vecinas y mide la estrella con PSF o con una apertura de $3\times3$~píxeles sobre la imagen
descontaminada. Las versiones \texttt{cal\_*} fijan la mediana al flujo que predice Gaia, lo que
corrige la dilución en escala absoluta. Con TGLC, \citet{Han2025} encuentran que los pipelines públicos
subestiman los radios planetarios en una mediana de 6~\%.

En MAST, TGLC solo existe para S5/S6 en estos objetos. El resto de los sectores se intentó correr
localmente (notebook 1c), pero el servicio TESScut estuvo caído (errores 502). Una salida inesperada:
\textbf{desde el sector 94, QLP usa la fotometría de TGLC} \citep{Petitpas2026}, así que las curvas
QLP de S98 son en la práctica TGLC de apertura $3\times3$. Así se usan aquí.

Atención: la serie TGLC mezcla versiones (HLSP v1 en S6, la implementación de QLP en S98 y tglc~0.8.1
local). Entre versiones hay diferencias de 10 a 25~\% en $D_1$ (en TIC~138829403 S6, por ejemplo, 0.055
en el HLSP y 0.050 local). Por eso $D_1$ de TGLC se usa como referencia de la dilución, pero no para
seguir $D_1$ de un sector a otro.

\subsection{QLP, CDIPS y PDCSAP}
\begin{itemize}
\item \textbf{QLP} \citep{Huang2020,Kunimoto2021}: fotometría de apertura del MIT (FIPHOT), sin
  corrección de dilución, hasta S93.
\item \textbf{CDIPS} \citep{Bouma2019}: imagen diferencia pensada para cúmulos jóvenes; se usa la
  apertura 2 (\texttt{IRM2}). Tampoco corrige la dilución.
\item \textbf{PDCSAP} de SPOC (cadencia de 2~min) y TESS-SPOC (FFI) \citep{Jenkins2016,Caldwell2020}:
  apertura corregida por la fracción de contaminación del TIC (\texttt{CROWDSAP}). En la misión
  primaria (sectores 1--26) SPOC sobrestimaba el fondo en campos poblados, lo que \emph{profundiza} los
  eclipses; se corrigió a partir del sector 27 \citep{Burt2020}.
\end{itemize}

\begin{table}[H]
\centering
\caption{Sectores disponibles por fotometría para los 13 sistemas revisados.}
\label{tab:disponibles}
\small
\begin{tabular}{lllllll}
\toprule
Sistema & TIC & \tx & TGLC & QLP ($<$S94) & CDIPS & PDCSAP \\
\midrule
1  & 176775425 & 6, 33, 87, 98 & 6, 87, 98 & 6, 33, 87 & 6, 33 & 33 \\
3  & 11412481  & 6, 98 & 6, 98 & 6 & -- & -- \\
4  & 50897998  & 6, 32, 98 & 6, 98 & 6, 32 & -- & 6, 32, 98 \\
5  & 249074181 & 5, 32, 98 & 5 & -- & -- & -- \\
6  & 72921630  & 6, 98 & 6 & -- & -- & -- \\
7  & 264483228 & 6, 32, 98 & 6, 98 & 6, 32 & 6, 32 & 6, 32 \\
8  & 11300044  & 6, 98 & 6 & 6 & -- & -- \\
9  & 427353034 & 6, 32, 98 & 6, 98 & 6, 32 & 32 & -- \\
10 & 72829000  & 6, 98 & 6 & 6 & -- & -- \\
11 & 138829403 & 6, 32, 98 & 6, 98 & 6, 32 & 6, 32 & 6, 32 \\
12 & 264741359 & 6, 32, 98 & 6, 98 & 6, 32 & 6, 32 & -- \\
16 & 427395536 & 6, 32, 98 & 6, 98 & 6 & 32 & -- \\
18 & 427347964 & 6, 32, 98 & 6, 98 & 6, 32 & 6, 32 & -- \\
\bottomrule
\end{tabular}
\end{table}

% ============================================================
\section{Dilución: por qué cambia la profundidad sin que cambie la estrella}
\label{sec:dilucion}
% ============================================================
Si una fracción $c$ del flujo en la apertura no es de la binaria, el eclipse observado es
\begin{equation}
D_{\rm obs}=(1-c)\,D,
\label{eq:dilucion}
\end{equation}
con $D$ la profundidad real. La fracción $c$ es lo que en los ajustes se llama tercera luz, $\ell_3$.
Hay dos consecuencias:
\begin{enumerate}
\item $c$ depende de la fotometría y \textbf{del sector}. Los píxeles de TESS miden 21'' y la PSF se
  extiende más de un minuto de arco, así que en Orión siempre entra luz de vecinas y de la nebulosa. En
  cada sector la estrella cae en otra cámara, con otra orientación del campo y otra forma de la PSF; la
  apertura y el anillo de fondo se recalculan. El resultado es un $c$ distinto por sector, y con él un
  $D_{\rm obs}$ distinto.
\item El cociente de profundidades \textbf{no depende de $c$}:
\begin{equation}
\frac{D_{2,\rm obs}}{D_{1,\rm obs}}=\frac{(1-c)\,D_2}{(1-c)\,D_1}=\frac{D_2}{D_1}.
\end{equation}
  Un cambio de $\Dum$ entre sectores no puede ser dilución. Es la prueba más limpia de que algo
  cambió en las estrellas.
\end{enumerate}

La tabla~\ref{tab:dilucion} da $c=1-D_{\tx}/D_{\rm ref}$ para cada fotometría de referencia en el mismo
sector. Con TGLC, que resta las vecinas, $c$ va de 0.2 a 0.6 en la mayoría de los sistemas.

\begin{table}[H]
\centering
\caption{Fracción de flujo ajeno en las curvas de \tx, $c=1-D_{1,\tx}/D_{1,\rm ref}$, en el mismo
sector y a 30~min. Un $c$ negativo indica que la referencia es la más somera (por ejemplo, PDCSAP con el
fondo sobrestimado o una estrella saturada).}
\label{tab:dilucion}
\small
\begin{tabular}{llcccc}
\toprule
TIC & Sector & TGLC & CDIPS & QLP & PDCSAP \\
\midrule
176775425 & 6  & 0.49 & 0.44 & 0.31 & -- \\
11412481  & 6  & 0.47 & --   & 0.32 & -- \\
11300044  & 6  & 0.43 & --   & 0.32 & -- \\
427353034 & 6  & 0.48 & --   & 0.33 & -- \\
427395536 & 6  & 0.60 & --   & 0.38 & -- \\
264741359 & 6  & 0.37 & 0.24 & 0.17 & -- \\
72921630  & 6  & 0.34 & --   & --   & -- \\
249074181 & 5  & 0.30 & --   & --   & -- \\
138829403 & 32 & --   & 0.13 & 0.18 & 0.16 \\
427347964 & 6  & 0.18 & 0.15 & 0.05 & -- \\
264483228 & 32 & --   & $-0.01$ & 0.00 & 0.00 \\
50897998  & 32 & --   & --   & 0.19 & 0.01 \\
\bottomrule
\end{tabular}
\end{table}

Una comprobación externa: \citet{KounkelStassun2024} estudiaron 2MASS~J05454167$-$0004024, que es
nuestro sistema~6 (TIC~72921630), con eleanor en S6. Ajustaron una tercera luz de $\ell_3=0.31$ en la
banda de TESS y la atribuyen a una compañera terciaria. Nuestra comparación de \tx\ con TGLC en el mismo
sector da $c=0.34$. Su periodo es $\porb=1.817212$~d, mientras que \texttt{priors\_v2.csv} tiene
1.82071~d; la diferencia ($3.5\times10^{-3}$~d) basta para achatar el eclipse plegado y hay que revisarla.

% ============================================================
\section{Método}
\label{sec:metodo}
% ============================================================

\subsection{Profundidad de cada eclipse}
En lugar de una profundidad por sector se mide cada eclipse por separado. Con el flujo normalizado sin
eclipses (biweight de \texttt{wotan}, \citealt{Hippke2019}) y la línea base $B$ del sector, para cada evento
se busca el mínimo de la mediana móvil con una ventana de $0.3\,W$ (nunca menor que dos cadencias) y el
centro libre en $\pm W/8$ alrededor de la efeméride:
\begin{equation}
D_k=B-\min_{c\in[-W/8,\,W/8]}\ \operatorname{med}\{F_i:\ |\phi_i-c|<h\},\qquad
\sigma_{D_k}=\sqrt{\tfrac{\pi}{2}}\,\frac{\hat\sigma}{\sqrt{n_k}}.
\end{equation}
Como cada eclipse se mide en su propio mínimo, un periodo con un pequeño error ya no achata la
profundidad: con el estimador por sector, la deriva acumulada en S98 abría el eclipse plegado en varias
trazas y lo hacía más somero. La profundidad de referencia de un sector es la mediana de sus eclipses.
Se guarda también el nivel de brillo fuera de eclipse en el momento de cada eclipse, para buscar la firma
de las manchas.

\subsection{Cadencia común}
El orden temporal de los sectores coincide con el de sus cadencias: 30~min en S5/S6, 10~min en
S32/S33 y 3.3~min en S87/S98. Una exposición larga promedia el fondo del eclipse con sus flancos
\citep{Kipping2010}, así que un eclipse corto se ve más somero, y más si es en V que si tiene fondo plano.
Eso puede imitar una tendencia temporal, tanto en $D_1$ como en $\Dum$. Para quitarlo, todas las curvas
se promedian a 30~min antes de medir. En el sistema~1, por ejemplo, promediar S33 a 30~min baja $D_1$ de
0.220 a 0.208 con el estimador por sector y de 0.213 a 0.181 eclipse por eclipse (figura~\ref{fig:cadencia}).

\begin{figure}[H]
\centering
\includegraphics[width=0.95\textwidth]{fig4_cadencia.pdf}
\caption{Efecto de la cadencia en TIC~176775425 (eclipse de 1.3~h). Azul: cadencia nativa; rojo: la misma
curva promediada a 30~min, como S6.}
\label{fig:cadencia}
\end{figure}

\subsection{Comparación entre fotometrías}
Cada fotometría trata la contaminación de otra forma, así que sus profundidades absolutas no tienen por
qué coincidir. \citet{Baum2025} lo advierten: los cambios de profundidad solo se pueden seguir \emph{dentro}
de un mismo conjunto de datos. Por eso se compara el \textbf{patrón entre sectores},
$D_1(s)/\langle D_1\rangle$ dentro de cada fotometría, y el cociente $\Dum$, que es comparable entre
fotometrías. Para cada fotometría y sistema se calculan además:
\begin{itemize}
\item un ANOVA sobre los eclipses individuales: ¿difieren las medias de los sectores más de lo que se
  dispersan los eclipses dentro de cada sector? ($p_{\rm entre}$);
\item la correlación de Spearman entre $D_k$ y el brillo fuera de eclipse dentro de cada sector
  ($p_{\rm brillo}$). Si la estrella eclipsada tiene manchas no ocultadas, el eclipse cambia con el
  brillo.
\end{itemize}

% ============================================================
\section{Resultados por sistema}
\label{sec:resultados}
% ============================================================

\begin{figure}[H]
\centering
\includegraphics[width=\textwidth]{fig2_patron_30min.png}
\caption{Patrón entre sectores por fotometría, con todas las curvas a 30~min. Para cada sistema, a la
izquierda $D_1(s)$ dividido por el promedio de esa fotometría y a la derecha $\Dum$.}
\label{fig:patron}
\end{figure}

\paragraph{Variación real ($\Dum$ cambia en todas las fotometrías).}
\begin{itemize}
\item \textbf{Sistema 1 (TIC 176775425).} $\Dum$ sube de 0.534--0.551 en S6 a 0.616--0.624 en S98 en
  \tx, QLP, CDIPS y TGLC. $D_1$ no tiene un patrón común: \tx\ da S98 $+7$~\% y la serie TGLC/QLP lo
  da más somero. El secundario se hace más profundo respecto del primario en unos 12~\% en siete años.
\item \textbf{Sistema 9 (TIC 427353034).} $\Dum$ baja de $\sim0.39$ a 0.36 en \tx, QLP y TGLC. Por
  otro lado, el $D_1$ de \tx\ en S98 cae 24~\%, cosa que no ocurre en las otras: eso es sistemático de
  \tx\ (la contaminación de S98 es mayor).
\item \textbf{Sistema 12 (TIC 264741359).} $D_1$ estable en \tx\ ($\pm1.5$~\%) y en CDIPS; $\Dum$ baja de
  0.387 a 0.357 y 0.342 en \tx, QLP, CDIPS y TGLC. El secundario se hace más somero. La correlación con el
  brillo es significativa ($p_{\rm brillo}\sim10^{-6}$), así que hay manchas. Gaia tiene una compañera a
  1.3'' ($\Delta T\approx1.6$~mag); se podría verificar de cuál viene la señal con TESS\_localize
  \citep{HigginsBell2023}.
\item \textbf{Sistema 18 (TIC 427347964).} Todas las fotometrías ven S32 más profundo y S98 $\sim35$~\% más
  somero, y $\Dum$ sube en S98. Es el único sistema con un cambio real de $D_1$. Como $\Dum$ también
  cambia, no es solo una tercera luz que varía: algo cambia en las propias estrellas.
\end{itemize}

\paragraph{Sistemático de \tx.}
\begin{itemize}
\item \textbf{Sistema 11 (TIC 138829403).} \tx\ da S32 $7$~\% más somero que S6; CDIPS, PDCSAP y QLP dan
  S6 y S32 iguales dentro de $\pm2$~\%. $\Dum$ estable.
\item \textbf{Sistema 4 (TIC 50897998).} \tx\ y PDCSAP (cadencia de 2~min, independiente) dan S98
  $\sim5$~\% más somero, con $\Dum$ constante (0.582--0.588). Es una dilución común a la fotometría de
  apertura de una estrella saturada ($T=5.5$); TGLC y QLP~S98 no sirven aquí.
\end{itemize}

\paragraph{Consistentes.}
Los sistemas 3, 7, 8, 10 y 16 no muestran un salto entre sectores cuando se miden eclipse por eclipse y a
la misma cadencia. Lo que marcaba el 1b era el estimador por sector (efeméride y cadencia). En el sistema~7,
$\Dum$ sube levemente (0.068 a 0.074) en todas las fotometrías; puede ser real pero es pequeño.

\paragraph{Sin comparación independiente.}
Los sistemas 5 (TIC 249074181) y 6 (TIC 72921630) solo tienen curvas de \tx\ en S32/S98. Las dos muestran
saltos grandes en $D_1$ y en $\Dum$ (en el sistema~5, $\Dum$ va de 0.68 a 1.41), pero son estrellas débiles
($T=13.9$ y 14.8), con mucha contaminación ($c\approx0.3$). Hace falta TGLC local para decidir.

\begin{figure}[H]
\centering
\includegraphics[width=\textwidth]{fig3_plegadas_clave.pdf}
\caption{Eclipse primario plegado por sector (colores) en cada fotometría (columnas), para los cuatro
sistemas con variación real de $\Dum$. Las líneas discontinuas marcan $1-D_1$ de cada sector, medido eclipse
por eclipse.}
\label{fig:plegadas_clave}
\end{figure}

% ============================================================
\section{Discusión: la tercera luz en el modelado de las curvas}
\label{sec:tercera_luz}
% ============================================================

\subsection{El efecto en los parámetros del ajuste}
La sección~\ref{sec:dilucion} mostró que las curvas de \tx\ tienen entre 15 y 50~\% de flujo que no viene
de la binaria. En el modelado esa luz se llama \emph{tercera luz}: un flujo constante, que no se eclipsa,
y que entra a la apertura desde estrellas vecinas, la nebulosa o una tercera estrella ligada al sistema.
Si una fracción $c$ del flujo normalizado es tercera luz, el modelo observado es
\begin{equation}
F_{\rm obs}(t)=(1-c)\,\frac{F_{\rm bin}(t)}{F_{\rm bin,0}}+c,
\label{eq:tercera_luz}
\end{equation}
con $F_{\rm bin,0}$ el flujo de la binaria fuera de eclipse, de modo que $D_{\rm obs}=(1-c)\,D$.

Los ajustes con starry de este trabajo suponen $c=0$: las amplitudes de las dos estrellas suman uno con
priors estrechos. Para reproducir un eclipse diluido, el modelo compensa con la geometría, ya sea con una
secundaria más pequeña ($k=R_B/R_A$ menor) o con un eclipse más rasante (inclinación menor). Como
$R_2\propto\sqrt{D}$, una tercera luz de $c\approx0.3$ se traduce en un $R_2$ subestimado en $\sim20$~\%,
y con él la masa y la temperatura de la secundaria que se obtienen de las isócronas. Además, como $c$
cambia de un sector a otro, un modelo con $c=0$ convierte esas diferencias en geometrías distintas para cada
sector, como se vio en TIC~427353034 S98 y TIC~50897998 S98.

\subsection{Cómo se ha abordado en otros trabajos}
Los enfoques publicados se pueden agrupar en cuatro, y en la práctica suelen combinarse.

\paragraph{Tercera luz como parámetro libre, por sector.} Es el tratamiento más extendido con TESS.
En la serie \emph{Rediscussion of eclipsing binaries}, Southworth ajusta la tercera luz con JKTEBOP como
un parámetro más, por separado en cada sector o periodo de observación \citep[por ejemplo,][]{Southworth2024CWEri}.
En el ajuste de tránsitos con \texttt{juliet} \citep{Espinoza2019}, la dilución es un parámetro por instrumento
y puede dejarse libre por sector, aunque muchos parámetros de dilución libres dificultan la convergencia.
PHOEBE~2.2 \citep{Jones2020} incorpora la tercera luz como flujo o como fracción del total.

\paragraph{Información externa para romper la degeneración.} La tercera luz está correlacionada con la
inclinación y con $k$, sobre todo en eclipses parciales, así que suele acotarse con información que no viene
de la curva de luz: el cociente de flujos espectroscópico \citep{David2016}; ajustes en varias bandas, porque
$c$ depende del color de la contaminante y la geometría no \citep[en Orión,][ajustan $\ell_3$ por banda y obtienen
$\ell_3=0.31$ en TESS para nuestro sistema~6]{KounkelStassun2024}; o la dinámica de una terciaria, a partir de los
tiempos de eclipse y la velocidad sistémica \citep{GomezMaqueoChew2019}.

\paragraph{Quitar la contaminación en la fotometría.} Con fotometría PSF que modela y resta las estrellas de
Gaia, como TGLC \citep{HanBrandt2023}, la dilución se corrige antes del ajuste; \citet{Han2025} muestran que, sin
esto, los radios planetarios de TESS salen subestimados en un 6~\% en promedio. SPOC corrige con la fracción de
contaminación del TIC (\texttt{CROWDSAP}), que \citet{Baum2025} revisan sector por sector antes de interpretar
cambios de profundidad. En ajustes conjuntos con fotometría terrestre de alta resolución angular, la dilución de
ésta se fija en uno y solo se ajusta la de TESS.

\paragraph{Comparación de modelos.} Varios trabajos ajustan con y sin tercera luz y eligen con un criterio
estadístico (por ejemplo, el BIC), de modo que solo se incluye cuando los datos lo justifican.

\subsection{Aplicación a esta muestra}
No contamos con cocientes de flujo espectroscópicos ni con fotometría terrestre, pero sí con dos elementos que
rara vez se tienen juntos: varios sectores con contaminación distinta y una medición independiente de $c$ por
sector, la comparación con TGLC (tabla~\ref{tab:dilucion}). El tratamiento más directo sería ajustar una sola
geometría para todos los sectores de un sistema, con un $c$ propio en cada sector y un prior centrado en el
valor de TGLC (con $\sigma_c\approx0.1$, para cubrir las diferencias entre versiones de TGLC), y comprobarlo
con el ajuste de las curvas TGLC de S98 con $c\approx0$. Este esquema tiene además una ventaja diagnóstica:
un $c$ por sector no puede reproducir un cambio de $D_2/D_1$, de modo que en los sistemas 1, 9, 12 y 18 el
cambio de $D_2/D_1$ debería seguir apareciendo como residuo si es real.

Mientras no se incluya la tercera luz en el modelo, los radios y masas de las secundarias derivados aquí deben
tomarse como \textbf{límites inferiores}.

% ============================================================
\section{Discusión: la luminosidad de la primaria del sistema 11 (TIC 138829403)}
\label{sec:sistema11}
% ============================================================
El tipo espectral de la literatura para la primaria del sistema~11 es B3.5 (FAST; \texttt{Macro\_21aug2026.csv}),
que en la escala de Pecaut \& Mamajek corresponde a $T_{\rm eff}\approx16\,700$~K. Con la fotometría de Gaia y
2MASS, la paralaje corregida ($\varpi=2.336\pm0.025$~mas, $d\approx430$~pc, RUWE $=1.34$) y una extinción de
$A_V=1.3$ (MassAge, anclado en $J$), la luminosidad resulta $\log L/L_\odot=2.25$. A 10~Myr, la isócrona PARSEC
de secuencia principal predice para esa temperatura $\log L/L_\odot=2.62$ ($M\approx4.6\,M_\odot$,
$R\approx2.5\,R_\odot$): la estrella queda $0.37$~dex (un factor 2.3) \textbf{por debajo} de la secuencia principal,
en una zona del diagrama HR que ninguna isócrona ocupa. Por eso MassAge no encuentra solución en la fase
pre-secuencia principal y, en la fase completa, devuelve una edad sin sentido ($\sim$100--330~Myr).

Las explicaciones posibles son:
\begin{itemize}
\item \textbf{Extinción subestimada.} Para subir la luminosidad 0.37~dex hacen falta $\sim0.9$~mag más en la banda
  $J$, es decir $A_V\approx3$--4 en lugar de 1.3. MassAge estima $A_V$ comparando colores observados con los
  intrínsecos de la temperatura dada; si la $T_{\rm eff}$ es menor que la supuesta, la extinción derivada también
  se subestima.
\item \textbf{Tipo espectral o temperatura sobreestimados.} El TIC~8.2 da $T_{\rm eff}=8\,388$~K a partir de colores,
  muy distinto del B3.5 espectroscópico. Si la estrella fuera más fría, con la misma fotometría quedaría en otra
  posición del diagrama HR. La $T_{\rm eff}$ del TIC no está corregida por enrojecimiento y no es confiable aquí,
  pero la discrepancia sugiere revisar la clasificación espectral.
\item \textbf{Distancia.} Para explicar el déficit con la distancia, la paralaje tendría que ser $\sim1.5$ veces
  menor; con $\sigma_\varpi/\varpi\approx1$~\% y RUWE de 1.34 eso es poco probable.
\end{itemize}

Este problema afecta a los parámetros de la primaria y, a través de ellos, a los de la secundaria: el radio
$R_1=1.63\,R_\odot$ del catálogo (Stefan--Boltzmann) es menor que los 2.0--2.2~$R_\odot$ del ajuste con starry y que
los $\sim2.5\,R_\odot$ esperados para una B3.5, de modo que $R_2=R_1\sqrt{D_1}$ y la masa de la secundaria que se
deriva de las isócronas ($\sim0.06\,M_\odot$) están subestimados. No afecta, en cambio, a la edad asignada al
sistema, que proviene de sus hermanas cinemáticas ($10.8\pm1.8$~Myr) y no de la posición de la primaria en el
diagrama HR. Resolverlo requiere un espectro que confirme el tipo espectral y una estimación independiente de la
extinción (por ejemplo, a partir de mapas 3D de polvo o de los colores de estrellas vecinas a la misma distancia).



% ============================================================
\section{Implicaciones para el ajuste con starry}
% ============================================================
\begin{enumerate}
\item \textbf{Tercera luz por sector.} Las curvas de \tx\ deben ajustarse con un parámetro de dilución por
  sector, $D_{\rm obs}=(1-\ell_{3,s})\,D$ \citep[como en][]{KounkelStassun2024}. Los valores de la
  tabla~\ref{tab:dilucion} sirven como priors, con una incertidumbre de al menos 10~\% por las
  diferencias entre versiones de TGLC. Sin este parámetro, starry subestima el radio y el brillo
  superficial de la secundaria.
\item \textbf{Integración en la exposición.} Los sectores de 30~min deben modelarse integrando el flujo
  durante la exposición \citep{Kipping2010}; starry lo permite \citep{Luger2019}.
\item \textbf{Sistemas con $\Dum$ variable} (1, 9, 12 y 18). Un modelo único para todos los sectores no va a
  ajustar bien. Hay dos opciones: ajustar cada sector por separado, o dejar que el cociente de brillos
  superficiales (o un mapa de manchas) cambie entre sectores.
\item \textbf{Efemérides.} El estimador por sector es sensible al periodo; conviene una efeméride global con
  todos los sectores (siete años de base). El sistema~6 tiene un periodo publicado distinto del nuestro.
\end{enumerate}

% ============================================================
\section{Pendientes}
% ============================================================
\begin{itemize}
\item Terminar las curvas TGLC locales (notebook 1c), sobre todo para los sistemas 5 y 6; depende de que
  TESScut responda.
\item Comprobar el periodo del sistema~6 con el valor de \citet{KounkelStassun2024}.
\item Actualizar la regla automática del notebook 1e para que use las medidas a 30~min y separe la prueba de
  $D_1$ (sistemática o real) de la de $\Dum$ (cambio físico).
\item Localizar con TESS\_localize \citep{HigginsBell2023} la fuente de la variabilidad en el sistema~12.
\end{itemize}

% ============================================================
\section{Reproducibilidad}
% ============================================================
\begin{itemize}
\item \texttt{1c\_tglc\_extraccion.ipynb} (kernel \texttt{tglc}): TGLC local; los sectores de 200~s se
  promedian de a 3 imágenes antes del ajuste del ePSF.
\item \texttt{1d\_tglc\_profundidades.ipynb} (kernel \texttt{starry}): TGLC frente a \tx, eclipses
  individuales.
\item \texttt{1e\_profundidades\_todas\_fotometrias.ipynb} (kernel \texttt{starry}): las cinco fotometrías;
  usa las curvas de \texttt{Light\_Curve\_Files\_MAST/} y no depende de TESScut. Salidas en
  \texttt{resultados\_pipeline/profundidades\_v2/pipelines/}.
\item Números de esta nota: \texttt{numeros\_profundidades\_30min.csv}; veredicto automático del 1e:
  \texttt{veredicto\_automatico.csv}.
\end{itemize}

\bibliographystyle{plainnat}
\bibliography{referencias}

\end{document}
